\documentclass[10pt,a4paper]{amsart}
\usepackage[margin=19mm]{geometry}
\usepackage{amsmath,amssymb,mathtools}
\usepackage[scr=rsfs,cal=euler]{mathalfa}
\usepackage{xfrac}
\usepackage[unicode,hidelinks,bookmarks=false]{hyperref}
\newcommand{\ie}{i.e.\ }
\newcommand{\cf}{cf.\ }
\newcommand{\resp}{respectively\ }
\newcommand{\Subsection}{\subsection}
\providecommand{\nobreakdash}{\nobreak\hbox{-}}
\setlength{\emergencystretch}{2em}
\multlinegap0pt
\newcommand{\Ebb}{\mathbb E}
\newcommand{\Pbb}{\mathbb P}
\newcommand{\R}{\mathbb R}
\newcommand{\smin}{\sigma_{\min}}
\theoremstyle{plain}
\newtheorem{theorem}{Theorem}
\newtheorem{proposition}{Proposition}
\newtheorem{lemma}{Lemma}
\newtheorem{corollary}{Corollary}
\newtheorem{definition}{Definition}
\newtheorem{remark}{Remark}
\title[Gaussian Critical-Threshold Instability in Real Phase Retrieval]{Gaussian Critical-Threshold Instability in Real Phase Retrieval}
\author{Christian E. H\"aggblom, Christian E. H{\"a}ggblom}
\date{Source: arXiv:2609.28198v1 (CC BY 4.0)}
\begin{document}
\maketitle

\noindent\emph{Source and licence.} The delivered work is the article shown in the title above, version arXiv:2609.28198v1 (CC BY 4.0), distributed under the Creative Commons Attribution 4.0 International licence, as recorded in the arXiv licence metadata for this exact version and shown on the abstract page saved with this item. The full licence notice, which carries the licence URL and the address of that abstract page, is the bundled file \texttt{licenses/arxiv-2609.28198-CC-BY-4.0.txt}.
\par\smallskip

\noindent\textit{Editorial scope.} Counted unit: the article's Lemma (source label lem:conditional-linearization) (source0.tex lines 1559--1584, source label \texttt{lem:conditional-linearization}): ``Conditional linearization Uniformly for (s C ,M ), q C(t M) t M M,s (C) r M,s (C), where C r M,s (C) 2 o(t M 2), and M,s (C) : s M z 1 z (CC ) -1 z , z N(0,I m).''. It is delivered together with the article's complete original proof (source0.tex lines 1586--1878, one proof environment adjacent to the statement, 293 lines), copied line for line from the article's own text. Required context retained uncounted: lem:specialized-inverse-tail (lines 1047--1077); lem:inverse-tail-threshold-shift (lines 1367--1397); cor:central-inverse-operator-moments (lines 2934--2962); cor:row-conditioned-factor-moments (lines 3064--3096). Editorial formatting only: an AMS \texttt{amsart} preamble that restates the article's own macro definitions and its own theorem declarations, and the theorem environments the retained lines use; the article's own class file is neither shipped nor input; the retained environments are renumbered sequentially in order of appearance, whereas the article prints them with its own numbering; the article's equation numbering is not reproduced and every cross-reference inside the retained text resolves to the wrapper's numbering. No citation occurs inside any retained line, so the wrapper carries no bibliography; All retained bodies are exact copies of the bundled \texttt{sources/211560/source0.tex}, so no omission receipt applies. No asset-inclusion command occurs inside any retained span and the package ships no image asset.


\section*{Retained context: lem:specialized-inverse-tail (source0.tex lines 1047--1077)}

\begin{lemma}[Uniform specialized inverse-tail asymptotic]
\label{lem:specialized-inverse-tail}
Let \(X_s\in\R^{s\times s}\) be standard Gaussian, and let
\(K_s\in\R^{s\times n_s}\) be independent of \(X_s\). Assume that
\[
        K_sK_s^\top\succeq I_s
\]
and that, for every fixed \(0<p<\infty\), there is \(s_0(p)\) such that
\begin{equation}
        \Ebb\|K_s\|^p\le C_p,\qquad s\ge s_0(p)
        \label{eq:K-moment-bound}
\end{equation}
with constants independent of \(s\). All conclusions below are as \(s\to\infty\).
Let \(R_s\to\infty\) satisfy
\[
        R_s\ge e^{cs}
\]
for some fixed \(c>0\). Then
\begin{equation}
        \Ebb_{K_s}
        \sup_{r\in[R_s/2,\,2R_s]}
        \left|
        \Psi_s(K_s;r)
        -
        \frac{\gamma_s}{r}a_s(K_s)
        \right|^2
        =
        o\left(\frac{s}{R_s^2}\right).
        \label{eq:uniform-specialized-inverse-tail}
\end{equation}
\end{lemma}


\section*{Retained context: lem:inverse-tail-threshold-shift (source0.tex lines 1367--1397)}

\begin{lemma}[Quantitative stability under an additive threshold shift]
\label{lem:inverse-tail-threshold-shift}
Assume the hypotheses of Lemma~\ref{lem:specialized-inverse-tail}.  Let \(D_s\ge0\) be
jointly distributed with \(K_s\), independent of \(X_s\), and suppose that there is a
constant \(A<\infty\) such that, for every fixed \(0<p<\infty\) and all sufficiently large \(s\),
\begin{equation}
        \Ebb D_s^p\le C_p s^{Ap}.
        \label{eq:threshold-shift-moments}
\end{equation}
Let \(\theta_s\) be any random variable jointly distributed with \((K_s,D_s)\), independent
of \(X_s\), and satisfying \(|\theta_s|\le1\).  Set
\[
        \widetilde R_s
        :=
        R_s+\theta_sD_s.
\]
Then
\begin{equation}
        \Ebb_{K_s,D_s,\theta_s}
        \left|
        \Psi_s(K_s;\widetilde R_s)
        -
        \Psi_s(K_s;R_s)
        \right|^2
        =
        o\left(\frac{s}{R_s^2}\right).
        \label{eq:quantitative-threshold-shift}
\end{equation}
Here, on the negligible event \(\widetilde R_s\le0\), the expression
\(\Psi_s(K_s;\widetilde R_s)\) is interpreted as \(1\).
\end{lemma}


\section*{Retained context: cor:central-inverse-operator-moments (source0.tex lines 2934--2962)}

\begin{corollary}[Central inverse-operator moments]
\label{cor:central-inverse-operator-moments}
Fix \(\eta\in(0,1/2)\).  Let
\[
        C\in\R^{m\times M}
\]
be standard Gaussian, and set
\[
        s:=M-m.
\]
Assume that
\[
        \frac{s}{M}\in[\eta,1-\eta].
\]
Then, for every fixed \(0<p<\infty\), there are constants
\(M_0=M_0(\eta,p)\) and \(C_{\eta,p}<\infty\) such that, for
\(M\ge M_0\),
\begin{equation}
        \Ebb
        \left\|(CC^\top)^{-1/2}\right\|^p
        =
        \Ebb\,\sigma_m(C)^{-p}
        \le
        C_{\eta,p}M^{-p/2}.
        \label{eq:central-inverse-operator-moments}
\end{equation}
More generally, the estimate holds whenever \(s+1>p\), with a constant that may
depend on \(p\) and on a lower bound for \(s/M\).
\end{corollary}


\section*{Retained context: cor:row-conditioned-factor-moments (source0.tex lines 3064--3096)}

\begin{corollary}[Moments of the row-conditioned factor]
\label{cor:row-conditioned-factor-moments}
Under the assumptions of
Corollary~\ref{cor:central-inverse-operator-moments}, let
\(S\in\R^{m\times m}\) have the same singular values as \(C\), and let
\[
        G_1\in\R^{s\times m}
\]
be an independent standard Gaussian matrix.  Then, for every fixed
\(0<p<\infty\) and all sufficiently large \(M\),
\begin{equation}
        \Ebb\|S^{-1}\|^p
        \le
        C_{\eta,p}M^{-p/2},
        \label{eq:central-S-inverse-moments}
\end{equation}
and
\begin{equation}
        \Ebb\|G_1S^{-1}\|^p
        \le
        C_{\eta,p}.
        \label{eq:central-conditioned-factor-moments}
\end{equation}
Consequently, for
\[
        K:=(-G_1S^{-1},I_s),
\]
one has
\begin{equation}
        \Ebb\|K\|^p\le C_{\eta,p}.
        \label{eq:central-K-moments}
\end{equation}
\end{corollary}


\section{The counted unit: Lemma (source label lem:conditional-linearization) and its complete original proof}

\begin{lemma}[Conditional linearization]
\label{lem:conditional-linearization}
Uniformly for \(s\in\mathcal C_{\eta,M}\),
\begin{equation}
        q_C(t_M)
        =
        t_M\Theta_{M,s}(C)+r_{M,s}(C),
        \label{eq:conditional-linearization}
\end{equation}
where
\begin{equation}
        \Ebb_C|r_{M,s}(C)|^2=o(t_M^2),
        \label{eq:conditional-remainder-L2}
\end{equation}
and
\begin{equation}
        \Theta_{M,s}(C)
        :=
        \frac{\gamma_s}{\sqrt M}
        \Ebb_z
        \sqrt{1+z^\top(CC^\top)^{-1}z},
        \qquad
        z\sim N(0,I_m).
        \label{eq:theta-def}
\end{equation}
\end{lemma}

\begin{proof}
By rotational invariance, choose \(Q\in O(M)\), measurably as a function of \(C\), such
that
\[
        CQ=(S,0),
\]
where \(S\in\R^{m\times m}\) is invertible almost surely.  In particular,
\[
        CC^\top=SS^\top.
\]
Conditionally on \(C\), the matrix \(GQ\) is again standard Gaussian and is independent of
\(C\).  Write
\[
        GQ=(G_1,G_2),
\]
where
\[
        G_1\in\R^{s\times m},
        \qquad
        G_2\in\R^{s\times s}
\]
are independent standard Gaussian blocks.  Thus
\[
        \smin\begin{pmatrix}C\\G\end{pmatrix}
        \stackrel{d}{=}
        \smin L,
        \qquad
        L:=
        \begin{pmatrix}
        S&0\\
        G_1&G_2
        \end{pmatrix}.
\]

The matrices \(S\) and \(G_2\) are invertible almost surely, and
\[
        L^{-1}
        =
        \begin{pmatrix}
        S^{-1}&0\\
        -G_2^{-1}G_1S^{-1}&G_2^{-1}
        \end{pmatrix}.
\]
Set
\[
        K:=(-G_1S^{-1},I_s)\in\R^{s\times M},
        \qquad
        D:=\|S^{-1}\|.
\]
The lower block of \(L^{-1}\) is \(G_2^{-1}K\).  Consequently,
\begin{equation}
        \|G_2^{-1}K\|
        \le
        \|L^{-1}\|
        \le
        \|G_2^{-1}K\|+D.
        \label{eq:block-inverse-norm-equivalence}
\end{equation}

Let
\[
        R_M:=\frac{\sqrt M}{t_M}.
\]
Since
\[
        \left\{\smin(L)\le\frac{t_M}{\sqrt M}\right\}
        =
        \{\|L^{-1}\|\ge R_M\},
\]
we may write
\[
        q_C(t_M)
        =
        \Ebb_{G_1}
        \Pbb_{G_2}\{\|L^{-1}\|\ge R_M\mid C,G_1\}.
\]

For fixed \(C\) and \(G_1\), define
\[
        \Psi_{C,G_1}(r)
        :=
        \Pbb_{G_2}\{\|G_2^{-1}K\|\ge r\}.
\]
The norm comparison \eqref{eq:block-inverse-norm-equivalence} gives the pointwise
probability sandwich
\begin{equation}
        \Psi_{C,G_1}(R_M)
        \le
        \Pbb_{G_2}\{\|L^{-1}\|\ge R_M\mid C,G_1\}
        \le
        \Psi_{C,G_1}(R_M-D).
        \label{eq:block-inverse-probability-sandwich}
\end{equation}
On the negligible event \(R_M-D\le0\), the right-hand side is interpreted as \(1\).

We now verify the hypotheses of the inverse-tail results from the preceding subsection.
Since \(s\in\mathcal C_{\eta,M}\), one has \(s\asymp_\eta M\).  Moreover,
\[
        t_M=\frac{w_M}{K_M},
        \qquad
        w_M=e^{o(M)},
\]
so there is a constant \(c_\eta>0\) such that
\[
        R_M\ge e^{c_\eta s}
\]
for all sufficiently large \(M\), uniformly over
\(s\in\mathcal C_{\eta,M}\).

The matrix \(K\) is independent of \(G_2\) and satisfies
\[
        KK^\top\succeq I_s.
\]

Because \(S\) has the same singular values as \(C\),
Corollary~\ref{cor:central-inverse-operator-moments} gives, for every fixed
\(0<p<\infty\),
\begin{equation}
        \Ebb_C D^p
        =
        \Ebb_C\|S^{-1}\|^p
        =
        O_{\eta,p}(M^{-p/2}).
        \label{eq:D-central-moments}
\end{equation}
Moreover, Corollary~\ref{cor:row-conditioned-factor-moments} gives
\begin{equation}
        \Ebb_{C,G_1}\|K\|^p
        \le
        C_{\eta,p}.
        \label{eq:K-central-moments}
\end{equation}

All moment bounds here are needed only for fixed orders and sufficiently
large \(M\), uniformly over central \(s\). The estimates in the inverse-tail
proof depend only on these uniform constants and the exponential lower bound
for \(R_M\); their remainders are therefore uniform over central overlaps.
Thus the additive-threshold estimate
Lemma~\ref{lem:inverse-tail-threshold-shift}, applied jointly to the randomness
\((C,G_1)\) with \(X_s=G_2\), \(D_s=D\), and \(\theta_s=-1\), yields
\begin{equation}
        \Ebb_{C,G_1}
        \left|
        \Psi_{C,G_1}(R_M-D)
        -
        \Psi_{C,G_1}(R_M)
        \right|^2
        =
        o\left(\frac{s}{R_M^2}\right).
        \label{eq:block-shift-L2}
\end{equation}
Because \(s\asymp_\eta M\) and \(R_M=\sqrt M/t_M\),
\[
        \frac{s}{R_M^2}
        =
        \frac{s}{M}t_M^2
        \asymp_\eta t_M^2.
\]
Hence the right-hand side of \eqref{eq:block-shift-L2} is \(o(t_M^2)\).

Define
\[
        \widehat q_C
        :=
        \Ebb_{G_1}\Psi_{C,G_1}(R_M).
\]
By the sandwich \eqref{eq:block-inverse-probability-sandwich},
\[
        0
        \le
        q_C(t_M)-\widehat q_C
        \le
        \Ebb_{G_1}
        \left[
        \Psi_{C,G_1}(R_M-D)
        -
        \Psi_{C,G_1}(R_M)
        \right].
\]
Jensen's inequality and \eqref{eq:block-shift-L2} therefore imply
\begin{equation}
        \Ebb_C
        |q_C(t_M)-\widehat q_C|^2
        =
        o(t_M^2).
        \label{eq:qC-block-reduction-error}
\end{equation}

It remains to linearize \(\widehat q_C\).  Apply
Lemma~\ref{lem:specialized-inverse-tail} jointly to the random matrix
\[
        K=K(C,G_1),
\]
which is independent of \(G_2\).  The joint moment estimate
\eqref{eq:K-central-moments} is the relevant hypothesis; no uniform conditional
moment bound for every fixed \(C\) is needed.  We obtain
\begin{equation}
        \Psi_{C,G_1}(R_M)
        =
        \frac{\gamma_s}{R_M}
        \Ebb_u\|u^\top K\|_2
        +
        \xi_{M,s}(C,G_1),
        \label{eq:conditional-inverse-tail-expansion}
\end{equation}
where
\begin{equation}
        \Ebb_{C,G_1}
        |\xi_{M,s}(C,G_1)|^2
        =
        o\left(\frac{s}{R_M^2}\right)
        =
        o(t_M^2).
        \label{eq:conditional-inverse-tail-error}
\end{equation}

Averaging \eqref{eq:conditional-inverse-tail-expansion} over \(G_1\) gives
\[
        \widehat q_C
        =
        \frac{\gamma_s}{R_M}
        \Ebb_{G_1,u}\|u^\top K\|_2
        +
        \overline\xi_{M,s}(C),
\]
where
\[
        \overline\xi_{M,s}(C)
        :=
        \Ebb_{G_1}[\xi_{M,s}(C,G_1)\mid C].
\]
By Jensen's inequality and \eqref{eq:conditional-inverse-tail-error},
\begin{equation}
        \Ebb_C|\overline\xi_{M,s}(C)|^2
        =
        o(t_M^2).
        \label{eq:averaged-inverse-tail-error}
\end{equation}

Since \(\|u\|_2=1\),
\[
        \|u^\top K\|_2
        =
        \sqrt{1+\|u^\top G_1S^{-1}\|_2^2}.
\]
Also,
\[
        \frac{1}{R_M}
        =
        \frac{t_M}{\sqrt M}.
\]
It follows from \eqref{eq:qC-block-reduction-error} and
\eqref{eq:averaged-inverse-tail-error} that
\begin{equation}
        q_C(t_M)
        =
        \frac{\gamma_s t_M}{\sqrt M}
        \Ebb_{G_1,u}
        \sqrt{1+\|u^\top G_1S^{-1}\|_2^2}
        +
        r_{M,s}(C),
        \label{eq:qC-before-Gaussian-average}
\end{equation}
with
\[
        \Ebb_C|r_{M,s}(C)|^2=o(t_M^2).
\]

For fixed \(u\in S^{s-1}\), the row vector \(u^\top G_1\) is distributed as
\(z^\top\), where \(z\sim N(0,I_m)\).  Hence
\[
        \|u^\top G_1S^{-1}\|_2^2
        \stackrel{d}{=}
        z^\top S^{-1}S^{-\top}z.
\]
The matrices
\[
        S^{-1}S^{-\top}
        \qquad\text{and}\qquad
        (SS^\top)^{-1}=S^{-\top}S^{-1}
\]
have the same eigenvalues.  Since \(z\) is rotationally invariant and
\(SS^\top=CC^\top\),
\[
        \Ebb_{G_1,u}
        \sqrt{1+\|u^\top G_1S^{-1}\|_2^2}
        =
        \Ebb_z
        \sqrt{1+z^\top(CC^\top)^{-1}z}.
\]
Substituting this identity into \eqref{eq:qC-before-Gaussian-average} proves
\eqref{eq:conditional-linearization}--\eqref{eq:theta-def}.
\end{proof}

\end{document}
